\chapter{Conclusion \& Future Work}

\section{What We Discussed}

\begin{itemize}
  \item \textbf{Conventional AES architecture.}  
        We examined the standard AES--128, AES--192, and AES--256 architectures, including
        SubBytes, ShiftRows, MixColumns, and AddRoundKey operations.  
        Simulation results demonstrated full functional correctness across all three key sizes.
  
  \item \textbf{Baseline optimized architecture (Virtex-7 reference).}  
        We analyzed the architecture proposed in the reference paper, which combines the
        S-box and inverse S-box into a shared hardware structure and merges MixColumns with
        AddRoundKey operations.  
        This significantly reduces redundant hardware and provides a strong baseline for
        area-efficient AES design.

  \item \textbf{Clock-gating motivation.}  
        We reviewed literature on low-power AES implementations, observing that excessive
        switching activity in S-box, MixColumns, and KeySchedule units contributes heavily to
        dynamic power consumption.  
        This motivates applying selective clock gating on top of the optimized architecture to
        further reduce power and improve overall VLSI efficiency.

  \item \textbf{Methodology.}  
        A structured, three-stage methodology was followed:  
        (i) implement and verify the conventional AES architecture,  
        (ii) implement and analyze the Virtex-7 optimized architecture,  
        (iii) extend the design using clock gating and compare all three architectures in terms
        of resource utilization.

  \item \textbf{Verification and synthesis.}
        Waveform inspection, RTL schematics, and NIST testvector simulations confirmed
        correctness.
        Synthesis results established the baseline LUT, MUX, and register usage for the
        forthcoming comparison studies.

  \item \textbf{Pipeline architecture baseline.}
        The full-pipeline AES architecture was implemented and synthesized at baseline,
        establishing 38,766 LUTs and 2,763 flip-flops for AES-128. This architecture
        provides a high-area, high-throughput reference point for the optimization study.

  \item \textbf{FSM architecture baseline.}
        The FSM AES architecture was implemented and synthesized at baseline, establishing
        6,794 LUTs and 412 flip-flops for AES-128 --- an 82 percent reduction in LUT usage
        relative to the pipeline baseline through sequential hardware reuse.
\end{itemize}


\section{What We Built (to Date)}

\begin{itemize}
  \item \textbf{Fully functional conventional AES cores} for 128, 192, and 256-bit keys with
        verified encryption and decryption operations.
  \item \textbf{Comprehensive functional verification suite} using testbenches, waveform tracing,
        and correctness validation via NIST vectors.
  \item \textbf{Synthesis and resource analysis} establishing area and logic consumption of the
        implemented designs.
  \item \textbf{Dual-architecture SystemVerilog RTL} --- a full-pipeline implementation
        (pipeline top-level module) and an FSM-based implementation (FSM top-level module)
        supporting AES-128, AES-192, and AES-256, with a shared testbench controlled by
        a pipeline enable parameter.
  \item \textbf{Baseline synthesis results} for both architectures at AES-128, establishing
        the before-optimization LUT and flip-flop counts used as the comparison reference.
\end{itemize}


\section{How We Will Build the Rest}

\begin{enumerate}
  \item \textbf{Phase~1 --- xtime arithmetic replacement (active).}
        Replace EXP3/LN3 GF multiplication with xtime shift-and-XOR arithmetic in
        the MixColumns and InvMixColumns modules, and remove the 512-entry lookup
        tables from the lookup-table array module. This is the highest-impact single change and
        targets a 25--35 percent LUT reduction in the pipeline architecture and 25--30
        percent in the FSM architecture.

  \item \textbf{Phase~2 --- Clock enable gating on FSM registers (planned).}
        Apply CE guards to the packed register structs in the three FSM state modules to
        suppress unnecessary register toggling when the FSM is idle. The guard condition
        includes \texttt{r.ready == 1} as a third term for the cipher and inverse cipher FSMs
        to prevent the packed struct freeze from breaking the ready-clear cycle. This phase
        targets dynamic power reduction without altering functional behavior.

  \item \textbf{Final comparative synthesis.}
        Run Vivado synthesis for both pipeline and FSM architectures after each phase and
        compare LUT, flip-flop, and estimated power figures across: baseline, post-Phase~1,
        and post-Phase~2.
\end{enumerate}


\section{Limitations}

\begin{itemize}
  \item The current work focuses primarily on functional correctness and structural
        optimization; power figures are estimate-based until clock gating is fully implemented.
  \item Only a subset of low-power techniques has been explored; other methods such as data
        gating or operand isolation remain unexplored.
  \item The Virtex-7 optimized architecture preserves serial datapath assumptions and may not
        generalize to high-throughput or multi-core AES accelerators.
  \item Synthesis is performed on a single FPGA family; ASIC-specific optimizations are not yet
        analyzed.
\end{itemize}


\section{Future Work}

\begin{itemize}
  \item \textbf{Phase~1 synthesis completion.}
        After xtime replacement is confirmed correct via simulation, run full Vivado
        synthesis for AES-128, AES-192, and AES-256 on both architectures to capture
        actual (not estimated) post-Phase~1 LUT and flip-flop counts.

  \item \textbf{Phase~2 CE gating integration.}
        Apply and synthesize the clock enable guards on the FSM state registers.
        Measure the reduction in dynamic switching power relative to the ungated FSM baseline.

  \item \textbf{Cross-key-size comparison.}
        Extend the baseline and post-optimization synthesis to AES-192 and AES-256 and
        produce a complete comparison table covering all three key sizes for both architectures.

  \item \textbf{Composite-field S-Box (CFA) optimization.}
        Replace the lookup-table S-Box with a composite-field arithmetic implementation to
        further reduce the SubBytes LUT contribution --- deferred from this work due to the
        larger scope of the change.

  \item \textbf{ASIC generalization.}
        Evaluate whether the xtime and CE gating optimizations transfer to an ASIC flow
        and quantify area and power improvements in a standard-cell technology target.
\end{itemize}

\noindent
In summary, this work establishes verified baseline implementations of both a pipeline and
an FSM AES architecture, quantifies their area tradeoffs, and defines a two-phase
optimization roadmap. Phase~1 eliminates log-table GF multiplication overhead and Phase~2
suppresses idle-state switching through CE gating, together providing a systematic and
measurable path toward a resource-efficient, energy-aware AES hardware implementation.
